\section{Distributed Key Generation and FROST Signatures}

\subsection{$\pedpop$: Key Generation for FROST}

In addition to presenting the FROST threshold signature scheme,
this chapter presents
a simplified two-round DKG that builds on Schnorr signatures and Pedersen's VSS.
In later work,
this DKG was given the name $\pedpop$,
which I will also use to refer to this scheme.

$\pedpop$ is a two-round distributed key generation scheme shown in Figure~\ref{fg:pedpop},
Each participant performs the role of a dealer,
by generating secret shares and a commitment via Pedersen's VSS
with respect to a randomly sampled secret.
Then,
each participant sends the shares to their respective participant (saving one
share for itself),
and sends the commitment to all participants.
Additionally,
each participant sends a proof of knowledge (specifically, a
Schnorr signature) that demonstrates knowledge of the secret that is shared. In
doing so, $\pedpop$ protects against rogue key attacks~\cite{RistenpartY07}.
Without this proof of knowledge,
a rushing adversary (i.e., one that can speak last) could perform a rogue key
attack by picking  heir public key with respect to any other participant's public key,
effectively cancelling out key material or forcing the group public key to hit some target
value.



\begin{figure*}[!ht]
  \begin{framed}

    \begin{flushleft}
    \textbf{Round 1}
\begin{description}
  \item[1.] Every participant $P_i$ samples $t$ random values
        $(a_{i0}, \dots, a_{i(t-1))}) \rgets \mathbb{Z}_q$, and uses these values as
        coefficients to define a degree $t-1$ polynomial $f_i(x) =
        \sum_{j=0}^{t-1} a_{ij} x^j$.
      \item[2.] Every $P_i$ computes a proof of knowledge to the
        corresponding secret $a_{i0}$ by calculating
        $\sigma_i = (R_i, \mu_i)$,
        such that
        $k \rgets \mathbb{Z}_q$, $R_i = g^k$,
        $c_i = H(i, \Phi, g^{a_{i0}}, R_i)$,
        $\mu_i = k + a_{i0}\cdot c_i$,
        with $\Phi$ being a context string to prevent
        replay attacks.
      \item[3.] Every participant $P_i$ computes a public commitment
        $\Vec{C_i} = \langle \phi_{i0}, \dots, \phi_{i(t-1)} \rangle$, where $\phi_{ij}
          = g^{a_{ij}}$, $0 \leq j \leq t-1$
      \item[4.] Every $P_i$ broadcasts $\Vec{C_i}, \sigma_i$ to all other
        participants.
      \item[5.] Upon receiving $\Vec{C_\ell}, \sigma_\ell$ from
        participants $1 \leq \ell \leq n, \ell \neq i$, participant $P_i$ verifies
        $\sigma_\ell = (R_\ell, \mu_\ell)$, aborting on failure, by
        checking
          \(R_\ell \stackrel{?}{=} g^{\mu_\ell} \cdot
           \phi_{\ell0}^{-c_\ell}\), where
          \(c_\ell = H\left(\ell, \Phi, \phi_{\ell 0}, R_\ell \right)\).

        Upon success, participants delete $\{ \sigma_\ell : 1 \le \ell \le n \} $.
\end{description}

    \smallskip

    \textbf{Round 2}
    \begin{description}
      \item[1.] Each $P_i$ securely sends to each other
      participant $P_\ell$
        a secret share $(\ell, f_i(\ell))$, deleting $f_i$ and each share
        afterward except for $(i, f_i(i))$, which they keep for themselves.
      \item[2.] Each $P_i$ verifies their shares by calculating:
          $g^{f_\ell(i)} \iseq \prod_{k=0}^{t-1} \phi_{\ell k}^{i^k}$, aborting if the check fails.

      \item[3.] Each $P_i$ calculates their long-lived private
        signing share by computing $s_i = \sum_{\ell=1}^n f_\ell(i)$, stores
        $s_i$ securely, and deletes each $f_\ell(i)$.
      \item[4.] Each $P_i$ calculates their public verification
        share $Y_i = g^{s_i}$, and the group's public key
        $Y = \prod_{j=1}^n \phi_{j0}$.
        Any participant can compute the public verification share of
        any other participant by calculating
        \( \displaystyle
        Y_i = \prod_{j=1}^n \prod_{k=0}^{t-1} \phi_{jk}^{i^k }
        \).

    \end{description}
    \end{flushleft}
  \end{framed}
    \caption{\textbf{KeyGen.} A distributed
    key generation (DKG) protocol $\pedpop$ that builds upon
    the DKG by Pedersen~\cite{Pedersen91}. Our variant includes a protection
    against rogue key attacks by requiring each participant to prove knowledge
    of their secret value commits, and requires aborting on misbehaviour.}
  \label{fg:pedpop}
\end{figure*}


\subsection{The FROST Signature Scheme}

\begin{figure*}[!htb]
	%\begin{linedgamespec}
	\begin{pchstack}[boxed,center,space=1em]
		\begin{pcvstack}[space=1em]
			\procedure{$\Setup(\usecp)$}{
				(\Gr, q, g) \leftarrow \GroupGen(\usecp) \\
				\text{\algcom Select two hash functions} \\
				\HashNon, \HashSig : \{0,1\}^* \rightarrow \Z_p \\
				\pp \gets ((\Gr, q, g), \HashNon, \HashSig) \\
				\pcreturn \pp
			}
			\procedure{${\KeyGen}(n, t)$}{
        \text{\algcom Shown in the centralized setting,} \\
        \text{\algcom but a DKG can likewise be used.} \\
        \sk \randpick \Z_p; \threshpk \gets g^{\sk} \\
        \{(i, \sk_i) \}_{i \in \set{n}} \randpick \tss.\IssueShares(\sk, n, t) \\
        \text{\algcom Shamir secret sharing of $\sk$} \\
        \pcfor i \in \set{n}~\aggpk_i \gets g^{\sk_i}\\
        \pcreturn (\threshpk, \{ \aggpk_i \}_{i \in \set{n}},  \{ \sk_i\}_{i \in \set{n}})
			}
			\procedure{$\Sign()$}{
        \text{\algcom Local signer has index $k$} \\
				d_k, e_k \randpick \Z_p^2 \\
        D_k \gets g^{d_k};\ E_k \gets g^{e_k};\\
				\state_k \gets (D_k, d_k, E_k, e_k) \\
        \rho_k \gets (D_k, E_k) \\
				\pcreturn (\rho_k, \state_k)
			}
		\end{pcvstack}
		\begin{pcvstack}[space=1em]
			\procedure{$\Sign'(\state_k, \sk_k,  m, \{ \rho_i \}_{i \in \S})$}{
        \text{\algcom $\Sign'$ must be called at most once per $\state_k$} \\
        \text{parse}~(  D_k, d_k, E_k, e_k )  \gets \state_k \\
        \text{parse}~ \{ D_i, E_i \}_{i \in S} \gets \{ \rho_i \}_{i \in \S}\\
        \pcreturn \bot\ \pcif (D_k, E_k) \not\in \{ D_i,E_i \}_{i \in S} \\
        \pcfor i \in \S \pcdo \\
        \pcind a_i \gets \HashNon(k, \threshpk, m, \{ (i, D_i, E_i)\}_{i \in \S}) \\
        \aggR \gets \prod_{i \in \S} D_i \cdot {E_i}^{a_i}\\
        c \gets \HashSig(\threshpk, m, \aggR) \\
				z_k \gets d_k + e_k a_k + c \lambda_k \sk_k\\
				\text{\algcom $\lagrange_k$ is the $k^{th}$ Lagrange coefficient for $\S$}\\
				\rho'_k \gets z_k \\
				\pcreturn \rho'_k
			}
			\procedure{$\Combine(\{(\rho_i, \rho_i')\}_{i \in \S})$}{
        \text{parse}~(D_i, E_i) \gets \rho_i, z_i \gets \rho'_i, i \in \S\\
        \aggR \gets \prod_{i \in \S} D_i\cdot{E_i}^{a_i};~z \gets \sum_{i \in \S} z_i \\
				\sigma \gets (\aggR, z)\\
        \pcreturn \sigma
			}
			\procedure{$\Verify(\threshpk, m, \sigma)$}{
				\text{parse}~(\aggR,z) \gets \sigma\\
				c \gets \HashSig(\threshpk, m, \aggR)\\
				\pcif \aggR \cdot \threshpk^c = g^z \\
				\pcind \pcreturn 1\\
				\pcreturn 0
			}
		\end{pcvstack}
	\end{pchstack}
	%\end{linedgamespec}
	\caption{The FROST threshold signature scheme.
  The public parameters $\pp$ are implicitly given as input to all algorithms and protocols.
  FROST assumes an external mechanism to choose the set $\S \subseteq \{1, \ldots, n \}$ of signers, where $t \leq \abs{\S} \leq n$.
  $\S$ is required to be ordered to ensure consistency.
  Verification is identical to (single-party) Schnorr signatures.}
\label{fig:frost}
\end{figure*}



